\documentclass[aspectratio=169]{beamer}
\input{header}
\coursecode{JKSSB}

\title{Excel Data Tools}
\subtitle{Lesson 27 --- Sort, Filter, Charts and PivotTables}
\author{JKSSB Computer Awareness}
\date{\today}

\begin{document}
\frame{\titlepage}

\begin{frame}{What Excel Data Tools Do}
  \begin{itemize}
    \item Data entered in a worksheet must be \key{organised}, \key{summarised},
      and \key{presented}.
    \item \key{Sort} and \key{filter} arrange and narrow the rows.
    \item \key{Tables}, \key{charts}, and \key{PivotTables} organise and
      summarise the data.
    \item \key{Printing}, \key{protection}, and \key{what-if} analysis prepare
      and safeguard the result.
  \end{itemize}
  \begin{block}{The one idea}
    These tools change the \emph{view} of the data; they do not change the data
    itself.
  \end{block}
\end{frame}

\begin{frame}{Sorting}
  \begin{itemize}
    \item \key{Sort} rearranges rows into a chosen order.
    \item \key{Ascending} order: A to Z, or smallest number to largest.
    \item \key{Descending} order: Z to A, or largest to smallest.
    \item Sorting works on a \key{column}, while the whole row moves with it.
  \end{itemize}
  \begin{block}{Data > Sort}
    The Sort dialog can order by \key{more than one level} --- for example, by
    Department, and then by Salary within each department.
  \end{block}
\end{frame}

\begin{frame}{Multi-Level and Custom Sort}
  \begin{itemize}
    \item \key{Add Level} in the Sort dialog sorts by a second and a third key.
    \item A \key{custom list} sorts by a defined order, such as days of the week
      or months.
    \item If the range has titles, tick \key{My data has headers} so the header
      row is not sorted as data.
    \item Excel can also sort by cell colour or font colour.
  \end{itemize}
  \begin{alertblock}{Trap alert}
    Sorting a \key{single column} on its own scrambles the records. Sort the
    whole table, or use an Excel table.
  \end{alertblock}
\end{frame}

\begin{frame}{Filtering with AutoFilter}
  \begin{itemize}
    \item \key{Filter} shows only the rows that meet a condition; the other rows
      are \key{temporarily hidden}, not deleted.
    \item Turn it on with \key{Data > Filter}, shortcut
      \key{Ctrl}$+$\key{Shift}$+$\key{L}.
    \item Drop-down arrows appear in the header row of each column.
    \item Clicking an arrow narrows the list to the chosen values.
  \end{itemize}
  \begin{block}{Hidden, not removed}
    Clearing the filter (\key{Clear Filter}) brings back every row.
  \end{block}
\end{frame}

\begin{frame}{Filter Criteria}
  \begin{itemize}
    \item \key{Value filter}: keep one or more ticks from the list.
    \item \key{Condition filter}: Text Filters, Number Filters, and Date Filters
      (for example, Greater Than 500, or Begins With ``A'').
    \item \key{Top 10} and \key{Above Average} are ready-made condition filters.
    \item \key{Advanced Filter} handles complex criteria and can copy the result
      to another location.
  \end{itemize}
\end{frame}

\begin{frame}{Excel Tables (Format as Table)}
  \begin{itemize}
    \item An \key{Excel table} is a formatted range with a name and a table
      style.
    \item Create one with \key{Insert > Table}, shortcut \key{Ctrl}$+$\key{T}.
    \item A table auto-expands as new rows are typed, and formulas auto-fill.
    \item It provides banded rows, header filter buttons, and an optional
      \key{Total Row}.
  \end{itemize}
  \begin{block}{Structured references}
    In a table, formulas use names such as \key{Table1[Sales]} instead of cell
    addresses like \key{B2}, so they stay correct as the table grows.
  \end{block}
\end{frame}

\begin{frame}{Table vs Plain Range}
  \begin{center}
    \begin{tabular}{lll}
      \toprule
      \textbf{Point} & \textbf{Plain range} & \textbf{Excel table} \\
      \midrule
      Name & None & Assigned name (Table1) \\
      Grows & Manually & Auto-expands \\
      Formulas & Cell addresses & Structured references \\
      Filters & Manual AutoFilter & Built-in header buttons \\
      \bottomrule
    \end{tabular}
  \end{center}
  \vspace{0.25cm}
  \muted{A table is not a PivotTable; a table is the source data that a
  PivotTable can summarise.}
\end{frame}

\begin{frame}{Charts: Common Types}
  \begin{itemize}
    \item \key{Column} or \key{bar} chart: compares values across categories.
    \item \key{Line} chart: shows a trend, usually over time.
    \item \key{Pie} or \key{doughnut} chart: shows parts of one whole.
    \item \key{Scatter (XY)} chart: shows the relationship between two numeric
      variables.
  \end{itemize}
  \begin{block}{Linked to data}
    A chart is \key{linked} to its source range; edit the data and the chart
    updates.
  \end{block}
\end{frame}

\begin{frame}{Chart Parts}
  \begin{itemize}
    \item \key{Chart area} is the whole chart; \key{plot area} is where the data
      itself is drawn.
    \item A \key{data series} is one set of values; the \key{legend} names each
      series.
    \item The \key{category axis} (X) holds labels; the \key{value axis} (Y)
      holds numbers.
    \item A chart can float on a worksheet or be moved to a separate
      \key{chart sheet}.
  \end{itemize}
  \muted{Data labels print the exact value on each data point.}
\end{frame}

\begin{frame}{PivotTables}
  \begin{itemize}
    \item A \key{PivotTable} summarises a large table of data in seconds.
    \item Fields are \key{dragged} into four areas: \key{Rows}, \key{Columns},
      \key{Values}, and \key{Filters}.
    \item The \key{Values} area aggregates by \key{Sum}, \key{Count},
      \key{Average}, and more.
    \item Change the field layout and the summary is \key{re-computed} at once
      --- that is the ``pivot''.
  \end{itemize}
  \begin{block}{Source unchanged}
    A PivotTable does not edit the source data; it is only a summary view of it.
  \end{block}
\end{frame}

\begin{frame}{PivotTable Areas and PivotCharts}
  \begin{center}
    \begin{tabular}{p{0.24\textwidth}p{0.62\textwidth}}
      \toprule
      \textbf{Area} & \textbf{What goes there} \\
      \midrule
      Rows & Field values listed down the left side \\
      Columns & Field values spread across the top \\
      Values & Numbers that are summarised (Sum, Count, Average) \\
      Filters & Fields used to filter the whole PivotTable \\
      \bottomrule
    \end{tabular}
  \end{center}
  \vspace{0.2cm}
  \muted{A \key{PivotChart} is a chart built on a PivotTable; the two are
  linked.}
\end{frame}

\begin{frame}{PivotTable vs Source Data}
  \begin{itemize}
    \item The PivotTable \key{summarises} the source; it never replaces it.
    \item If the source data changes, the PivotTable must be \key{refreshed}.
    \item Deleting the source table breaks the PivotTable.
    \item A PivotTable is \key{not} the same as an Excel table.
  \end{itemize}
  \begin{alertblock}{Trap alert}
    A PivotTable ``summarises data but is not the source table itself''. Do not
    confuse the summary with the data.
  \end{alertblock}
\end{frame}

\begin{frame}{Printing: Page Setup}
  \begin{itemize}
    \item Print from \key{File > Print}; adjust layout on the \key{Page Layout}
      tab.
    \item \key{Orientation}: \key{Portrait} (tall) or \key{Landscape} (wide).
    \item Set \key{paper size}, \key{margins}, and \key{headers/footers} here.
    \item Preview in \key{Print Preview} before printing a large sheet.
  \end{itemize}
  \begin{block}{Keep it on one page}
    \key{Scaling} options such as ``Fit Sheet on One Page'' shrink the printout
    to a chosen width or number of pages.
  \end{block}
\end{frame}

\begin{frame}{Printing: Area, Titles and Gridlines}
  \begin{itemize}
    \item \key{Print Area}: select a range and \key{Set Print Area} to print
      only that part.
    \item \key{Print Titles} repeats chosen \key{header rows} or columns on every
      printed page.
    \item \key{Gridlines} are \key{not printed by default}; enable
      \key{Print} gridlines under Sheet Options.
  \end{itemize}
  \muted{Page breaks can be inserted to control where a new page starts.}
\end{frame}

\begin{frame}{Protection}
  \begin{itemize}
    \item \key{Protect Sheet} (\key{Review > Protect Sheet}) stops cells from
      being edited; selected actions may still be allowed.
    \item \key{Protect Workbook} protects the \key{structure} --- adding,
      deleting, or renaming sheets.
    \item Cells are \key{Locked} by default, but locking takes effect only
      \key{after} the sheet is protected.
  \end{itemize}
  \begin{alertblock}{Not the same as encryption}
    Protection restricts editing inside Excel; it is not the same as a password
    that encrypts the file.
  \end{alertblock}
\end{frame}

\begin{frame}{What-If Analysis (Awareness)}
  \begin{itemize}
    \item \key{Goal Seek} finds the input value needed to reach a target result.
    \item \key{Scenario Manager} stores and compares several sets of input
      values.
    \item A \key{Data Table} shows how an output changes as one or two inputs
      change.
    \item These tools live under \key{Data > What-If Analysis}.
  \end{itemize}
  \muted{Exam depth: recognise each tool and its one-line purpose.}
\end{frame}

\begin{frame}{Near-Neighbour: COUNTIFS vs SUMPRODUCT}
  \begin{itemize}
    \item To \key{count rows meeting several conditions at once}, use
      \key{COUNTIFS(range1, crit1, range2, crit2, ...)}.
    \item \key{COUNTIFS} applies \key{AND} logic across all its
      range-and-criteria pairs.
    \item \key{SUMPRODUCT((range1>=x)*(range2>=y))} also counts
      multi-criterion rows, by multiplying TRUE/FALSE results.
    \item For \key{OR} logic, add two COUNTIF or COUNTIFS results.
  \end{itemize}
  \begin{alertblock}{Trap alert (TRAP-14, PYQ-15)}
    COUNTIFS takes \key{paired} range-and-criteria arguments; SUMPRODUCT works
    on \key{arrays} and adds the already-multiplied boolean results directly.
  \end{alertblock}
\end{frame}

\begin{frame}{Trap Alert: Data-Tools Facts to Keep Straight}
  \begin{itemize}
    \item Sort and filter change the \key{view}; a filter \key{hides} rows, it
      does not delete them.
    \item A PivotTable \key{summarises} the source and does not replace it; it
      must be \key{refreshed} after changes.
    \item A PivotTable is \key{not} an Excel table.
    \item \key{COUNTIFS} counts with AND across paired arguments; \key{SUMPRODUCT}
      does the same with boolean arrays. (TRAP-14)
    \item Locked cells are protected only \key{after} the sheet is protected.
    \item Gridlines are \key{not printed} unless Print gridlines is enabled.
  \end{itemize}
\end{frame}

\begin{frame}{Lesson 27: Recall}
  \begin{itemize}
    \item \key{Sort} arranges rows --- ascending, descending, multi-level, and
      custom.
    \item \key{Filter} (\key{Ctrl}$+$\key{Shift}$+$\key{L}) shows only matching
      rows; it hides, not deletes.
    \item An \key{Excel table} (\key{Ctrl}$+$\key{T}) auto-expands and uses
      structured references.
    \item Charts (column, line, pie, scatter) are \key{linked} to their data.
    \item A \key{PivotTable} summarises with Rows, Columns, Values, and Filters;
      a \key{PivotChart} is a chart on a PivotTable.
    \item Printing: print area, print titles, scaling, and gridline options.
    \item \key{Protect Sheet} vs \key{Protect Workbook}; locking applies after
      protection.
    \item \key{Goal Seek}, \key{Scenario Manager}, and \key{Data Table} for
      what-if analysis.
    \item \key{COUNTIFS} vs \key{SUMPRODUCT} for multi-criterion counting
      (TRAP-14).
  \end{itemize}
\end{frame}
\end{document}
